
\begin{abstract}

\noindent
\textsc{VALOR}\footnote{\url{https://valor-fit.github.io}} is a neutrino oscillation fitting group originally established within the T2K experiment in 2010. Since its inception, the VALOR group has led multiple iterations of T2K’s electron-(anti)neutrino appearance, muon-(anti)neutrino disappearance, and joint three-flavour oscillation analyses, contributing to nearly all published T2K oscillation results.
\smallskip

\noindent
Following a large-scale code refactoring of the T2K analysis infrastructure, a generic, CPU-efficient, and highly flexible oscillation analysis framework—the \emph{VALOR Software Development Kit (SDK)}—was developed from the original VALOR codebase. The SDK has enabled the streamlined deployment of VALOR-style oscillation analyses across a broad range of current and future neutrino experiments.
\smallskip

\noindent
At present, VALOR-based oscillation analyses are in use or development for T2K beam data, JUNO atmospheric neutrino studies, SBN sterile neutrino oscillation and BSM searches, and advanced sensitivity simulations for DUNE and Hyper-Kamiokande. These analyses are implemented as thin layers of experiment-specific configuration and inputs on top of the shared SDK infrastructure.
\smallskip

\noindent
In this contribution, we present an overview of the VALOR SDK’s architecture and toolkit, and describe how it supports the construction, validation, and execution of sophisticated oscillation analyses across experiments.

\end{abstract}

\begin{keyword}
VALOR, SDK, Neutrino, Oscillations, Data Analysis
\end{keyword}